\section{Protezione della copia e crittografia White-Box}
\subsection{Contesto storico e tecniche di attacco}
Il contesto storico è quello dei primi anni 2000 -- periodo in cui erano appena entrate in commercio console come Xbox, GameCube e PlayStation 2. \\
Come anticipato in introduzione, il modello \textit{MATE} assume in partenza che tutto ciò che si trova all'esterno del processore (RAM, bus di sistema, periferiche, \ldots) sia esposto e potenzialmente sotto il controllo dell'avversario.\\
Il principio è lo stesso che giustifica la crittografia White-Box: chi esegue il programma è al contempo chi può attaccarlo. \\
Il contributo di \textit{Weidong Shi, Hsien-Hsin S. Lee, Chenghuai Lu e Tao Zhang}, pur non essendo specifico dell'ambito videoludico, esamina 
\hyperref[fonte:copy_protection]{\textbf{Attacchi e analisi del rischio per i sistemi hardware a supporto della protezione dalla copia del software}}~[\ref{fonte:copy_protection}]\footnote{Titolo originale: \textit{Attacks and Risk Analysis for Hardware Supported Software Copy Protection Systems}} e cita tra le principali tecniche di attacco hardware per soverchiare le difese software:
\begin{itemize}
    \item \textbf{Trace Analysis}: l'attaccante raccoglie informazioni del software protetto tramite \textit{logging} ed analizzando le tracce d'esecuzione dello stesso. Alcune primitive sono riconoscibili anche dopo essere state cifrate (ad esempio branch condizionali, indici e diversi tipi di variabili).
    \item \textbf{Mod-chip}: con il termine \textit{mod-chip} gli autori si riferiscono a dispositivi designati a mettere alla luce i meccanismi di protezione di copia. Possono anche essere destinati all'appropriazione del segnale di un altro dispositivo o a lanciare attacchi di spoofing ad altri dispositivi.
    \item \textbf{Emulatori}: è complesso combattere contro attacchi di emulazione senza utilizzare la cifratura; tuttavia, non è possibile parlare di protezione della copia del software se il suddetto può essere eseguito su un emulatore senza autorizzazione.
\end{itemize}
Scendendo ulteriormente nel dettaglio, gli autori analizzano anche altri contesti ed alcuni dei relativi attacchi:
\begin{itemize}
    \item Analizzando le \textbf{lazy authentication}\footnote{Con `lazy' ci si riferisce alla situazione nella quale la sorgente dei dati è utilizzata come un operando ed il risultato della computazione ottenuto utilizzando dati non autenticati è in grado di modificare lo stato del processore prima che l'integrità di 
    tali dati sia confermata.} \textbf{e le counter-mode encryption}\footnote{La counter-mode encryption (\textit{CTR}), invece di cifrare il messaggio direttamente, si cifra una sequenza di valori numerici incrementali per produrre un flusso di bit pseudo-casuali, detto pad. Il messaggio viene poi combinato con questo pad tramite XOR, ottenendo il ciphertext. 
    Fonte: \hyperref[fonte:ctr]{\textit{Comments to NIST concerning AES Modes of Operations: CTR-Mode Encryption}}~[\ref{fonte:ctr}], Lipmaa et al.}, si può parlare ad esempio di \textbf{attacco ATT}, ovvero \textit{''alter then trace attack''}, il quale prevede che l'attaccante sia in grado di modificare una parte del software 
    (programma o dati) bit-per-bit, che le istruzioni alterate possano essere eseguite e che l'integrità di tali istruzioni/dati non venga verificata tempestivamente né in modo rigoroso istruzione per istruzione.\\
    Il rischio delle autenticazioni ''Lazy'' può essere ridotto utilizzando \textit{control flow/memory fetch address obfuscation} hardware, nonostante questo possa aumentare notevolmente l'overhead e la complessità dell'hardware stesso.
    \item Riguardo il \textbf{Software slice}, una tecnica che prevede di mettere in sicurezza il software proteggendone solo alcune porzioni, l'attaccante può trattare le porzioni di codice protetto come \textit{black box} e cercare di creare un codice diverso dall'originale ma che riesca ad emettere lo stesso output dallo stesso input -- 
    naturalmente si tratta di un problema computazionale complesso se applicato a qualsiasi programma, ma l'avversario potrebbe essere interessato nel corrompere anche solo una \textit{slice} di programma, e potrebbe riuscire a farlo in un tempo ragionevole. 
    \item Gli autori del paper analizzano poi gli \textbf{Active Information Flow Attack}. Il contesto prevede che alcuni sistemi utilizzino lo stesso schema di cifratura e la stessa chiave sia per proteggere le istruzioni che i dati dinamici, esponendo il controllo della protezione all'utente tramite istruzioni come \texttt{enter security}, 
    \texttt{exit security}, \texttt{secure load} e \texttt{secure store}\footnote{Fonte: Table 3, \hyperref[fonte:copy_protection]{\textit{Attacks and Risk Analysis for Hardware Supported Software Copy Protection Systems}}~[\ref{fonte:copy_protection}], Shi et al.}.\\ Vengono portati due esempi di attacco a questi sistemi: nel primo, l'attaccante può semplicemente rimpiazzare parti di codice non protetti con contenuto malevolo;
    dopo che il codice infetto viene criptato, l'attaccante può ridirezionare il puntatore verso tale codice per farlo eseguire in un contesto protetto. Nel secondo esempio, l'attaccante cerca di sfruttare le \textit{linked list} per ridirezionare l'ultimo puntatore della lista ad un elemento segreto per poterlo rivelare in uno spazio non protetto. Le soluzioni
    a questo tipo di attacchi prevedono anche l'impiego di \textit{program analysis} o di diverse tecniche di compilazione.
    \item In alcuni contesti potrebbe essere necessario utilizzare \textbf{cifrature semplici o chiavi corte}, per motivi di performance o di limitazioni hardware. Anche in questo contesto, gli autori presentano due semplici esempi di attacco, uno che mira alla derivazione di un \textit{integrity code} breve a protezione di una porzione di codice tramite \textit{brute force attack}, 
    mentre l'altro che sfrutta le caratteristiche cicliche delle variabili contatore, come il fatto che incrementano in modo regolare, per individuarle (anche se cifrate -- ma con una chiave corta) e ricavarne le locazioni in memoria.
\end{itemize}

\subsection{Case Study: la sicurezza di Xbox (OG)}
Il seguente case study è basato sull'analisi di \textit{Michael Steil} nel paper \hyperref[fonte:xbox]{\textbf{17 Mistakes Microsoft Made in the Xbox Security System}}~[\ref{fonte:xbox}].\\
L'originale Xbox, introdotta sul mercato da Microsoft nel 2001, era stata ideata per essere utilizzata senza software copiati, applicazioni non ufficiali e sistemi operativi alternativi; tutti i sistemi di sicurezza, pertanto, erano stati implementati per adempiere a questo scopo.\\
In termini di hardware, per ridurre le tempistiche di sviluppo e costruzione, Microsoft ha optato per creare la sua console a partire da componenti per PC \textit{off the shelf}, mentre per il software si è affidata alle tecnologie proprietarie di Windows e DirectX come basi della console.\\
Nonostante gli errori compiuti da Microsoft in questa console siano diversi, vogliamo concentrarci solo sulla concretizzazione di quanto analizzato fino ad ora in questo capitolo.
\subsubsection{Trace Analysis e \textit{bus sniffing}}
L'autore cita il lavoro di Andrew "bunnie" Huang\footnote{\href{https://www.bunniestudios.com/bunnie/proj/anatak/xboxmod.html}{\textit{Bunnie's adventures hacking the Xbox}}}, al tempo studente PhD al MIT, il quale una volta disassemblata la console ha ricavato moltissime informazioni dalla flash memory:
analizzandone il contenuto, trovò nei 512 byte superiori il codice di una vecchia versione della ROM rimasta per errore nel build finale -- uno spazio che invece sarebbe dovuto rimanere libero, dedicato ad una porzione di memoria denominata d'ora in avanti in \textit{secret ROM}, destinata al setup della console e dal codice di decriptazione della memoria flash. 
Bunnie scoprì che questo codice era un interpreter per le tabelle in memoria flash, ed in aggiunta una funzione di decriptazione che somigliava ad \textit{RC4}\footnote{\href{https://en.wikipedia.org/wiki/RC4}{RC4}: Uno tra i più famosi e diffusi \textbf{algoritmi di cifratura a flusso a chiave simmetrica}, utilizzato ampiamente in protocolli quali l'SSL ed il WEP.}.\\
Dal momento che, indipendentemente da quali operazioni compiesse su questa porzione di memoria, la console si avviava correttamente, capì che doveva necessariamente essere sovrascritta ad ogni avvio, e che quello era il luogo dove si trovava la \textit{secret ROM}.\\
Sfruttando le risorse del laboratorio hardware del MIT, Huang costruì un hardware dedicato per `sniffare' il bus \textit{HyperTransport}\footnote{\href{https://en.wikipedia.org/wiki/HyperTransport}{HyperTransport (HT)}: Conosciuto anche come \textit{Lightning Data Transport}, è una tecnologia destinata all'interconnessione tra processori di un computer.}, 
ritenuto inaccessibile da Microsoft, estraendo così la \textit{secret ROM} completa e la chiave \textit{RC4}.
\subsubsection{Mod-chip}
Dopo aver ottenuto le chiavi di crittografia, comparvero i mod-chip: alcuni rimpiazzavano interamente la memoria flash, mentre altri ne `patchavano' solo una parte, ed utilizzavano per la gran parte il chip originale. \\
Se l'Xbox veniva avviata senza alcuna flash memory installata, essa tentava di leggere dal bus LPC; questa scoperta ha portato allo sviluppo di chip che inducevano la console a `credere' di non avere alcuna memoria installata, ed il \textit{fallback} era un chip custom che permetteva di copiare interi giochi sull'hard disk della console e di lanciarli da lì.
\subsubsection{Emulatori e software non autorizzato}
L'idea originale di Microsoft era quella di escludere qualsiasi software che non provenisse dal supporto originale o che non fosse autorizzato da Microsoft. Il design prevedeva che il sistema di sicurezza rendesse impossibile installare Linux sulla console, come anche una qualsiasi versione di software senza licenza.\\
Da un lato, l'idea rendeva il sistema di sicurezza più semplice e riduceva i punti di attacco, ma dall'altro ha triplicato i potenziali attaccanti: nonostante la comunità Open Source, gli sviluppatori \textit{homebrew} e i \textit{cracker} avessero interessi molto diversi tra loro, in questo caso potevano unirsi e attaccare congiuntamente il sistema di sicurezza Xbox.
\subsubsection{Lazy authentication e verifica non rigorosa}
Il sistema di sicurezza Xbox prevedeva che la \textit{secret ROM} verificasse l'autenticità del secondo bootloader (\textit{2bl}) prima di eseguirlo: a questo scopo, Microsoft scelse di utilizzare \textit{RC4} sia come algoritmo di decriptazione che come funzione di hash, confrontando gli ultimi 32 bit decifrati con una costante nota.\\
Tuttavia, RC4 è un cifrario a flusso che decifra ogni byte in modo indipendente: modificare un byte del testo cifrato non compromette la decifratura dei byte successivi. La verifica degli ultimi 32 bit risultava quindi priva di qualsiasi valore -- non esisteva alcun hash reale. Il \textit{2bl} poteva essere modificato arbitrariamente senza che il sistema se ne accorgesse.
\subsubsection{Cifrature semplici e chiavi corte}
I vincoli di spazio della \textit{secret ROM} -- grande solo 512 byte -- hanno costretto Microsoft a scegliere algoritmi crittografici estremamente compatti. La scelta ricadde su \textit{RC4} con una chiave a 16 byte, archiviata nella stessa \textit{secret ROM}.\\
Una chiave così corta, una volta che il bus \textit{HyperTransport} fu sniffato da Huang, risultò immediatamente estraibile. Nella "versione 1.1"\footnote{Chiamata così dalla community di hacker, si trattava di un update della ROM in seguito alla scoperta di Bunnie.} Microsoft aggiunse \textit{TEA} (\textit{Tiny Encryption Algorithm}) come hash, scelto verosimilmente per le sue dimensioni ridotte; tuttavia, come documentato da Schneier in \textit{Applied Cryptography}, TEA non è sicuro se 
utilizzato come funzione di hash -- invertendo il 16-esimo ed il 31-esimo bit di una parola a 32 bit si ottiene lo stesso valore di hash, rendendo semplicissima la manipolazione del \textit{2bl}.

